%% toc: P2 17 17.8
\section{Proofs for the quantified subformula property}
\label{sec:annex-p02-w02b-p02-020-subformula-proof}

Throughout, \(\Lang\) is a first-order language with terms
\(\operatorname{Term}(\Lang)\) (\cref{sec:fol-syntax-semantics-01}) and formulas
\(\operatorname{Fm}(\Lang)\) (\cref{sec:fol-syntax-semantics-02}); \(\varphi[t/x]\) is
the total capture-avoiding substitution of
\cref{def:p02-fol-capture-avoiding-substitution}; and
\(\operatorname{SI}(\chi)\), immediate propositional subformulas, literal
subformulas and successors are as in
\cref{def:p02-w02b-p02-020-instance-set,def:p02-w02b-p02-020-successor}.
Formulas are compared as strings.  For a fixed \(\chi\), the set
\(\operatorname{SI}(\chi)\) exists: \(\operatorname{Fm}(\Lang)\) itself
contains \(\chi\) and is closed under both clauses of
\cref{def:p02-w02b-p02-020-instance-set}, and the intersection of all sets of
formulas with these two properties again has them.  The same remark applies to
the set of literal subformulas.

\begin{annexitem}{annex:p02-w02b-p02-020-closure}{thm:p02-w02b-p02-020-closure}{Instance sets, chains and transitivity}
\begin{proof}
\emph{Clause 1.}  Write \(\mathrel{\mathsf R}\) for the relation of
\cref{def:p02-fol-subformula-instance}, so that \(\psi\mathrel{\mathsf R}\chi\)
says that \(\psi\) is a first-order subformula instance of \(\chi\); its
clause~3 ranges over \(t\in\operatorname{Term}(\Lang)\), as recorded in
\cref{def:p02-w02b-p02-020-instance-set}.  Fix \(\chi\) and put
\(A_\chi=\{\psi:\psi\mathrel{\mathsf R}\chi\}\).  By clause~1 of that
definition \(\chi\in A_\chi\); by its clauses 2 and 3, \(A_\chi\) is closed
under passing to immediate propositional subformulas and under passing from
\(Qx\,\theta\) to \(\theta[t/x]\).  Since \(\operatorname{SI}(\chi)\) is the
least set with these properties, \(\operatorname{SI}(\chi)\subseteq A_\chi\).
Conversely, let \(\mathsf R'\) be the relation that holds of \((\psi,\chi)\)
exactly when \(\psi\in\operatorname{SI}(\chi)\).  It satisfies the three
clauses of \cref{def:p02-fol-subformula-instance}: \(\chi\in
\operatorname{SI}(\chi)\); if \(\varphi\in\operatorname{SI}(\chi)\) and
\(\psi\) is an immediate propositional subformula of \(\varphi\), then
\(\psi\in\operatorname{SI}(\chi)\) by clause~1 of
\cref{def:p02-w02b-p02-020-instance-set}; and if
\(Qx\,\theta\in\operatorname{SI}(\chi)\), then
\(\theta[t/x]\in\operatorname{SI}(\chi)\) for every term \(t\) by clause~2.
Since \(\mathsf R\) is the least relation satisfying the three clauses,
\(\mathsf R\subseteq\mathsf R'\), that is, \(A_\chi\subseteq
\operatorname{SI}(\chi)\).

\emph{Clause 2.}  Call a sequence as in the statement a \emph{chain} from
\(\chi\), and let \(C_\chi\) be the set of last members of chains from
\(\chi\).  The one-term chain gives \(\chi\in C_\chi\).  If \(\psi\in C_\chi\)
is the last member of a chain, and \(\psi'\) is an immediate propositional
subformula of \(\psi\), or \(\psi=Qx\,\theta\) and \(\psi'=\theta[t/x]\) with
\(t\in\operatorname{Term}(\Lang)\), then appending \(\psi'\) gives a chain
ending in \(\psi'\).  So \(C_\chi\) contains \(\chi\) and is closed under both
clauses, and \(\operatorname{SI}(\chi)\subseteq C_\chi\).  Conversely, let
\(\chi=\psi_0,\ldots,\psi_n\) be a chain.  By induction on \(i\le n\),
\(\psi_i\in\operatorname{SI}(\chi)\): \(\psi_0=\chi\), and \(\psi_{i+1}\) is
obtained from \(\psi_i\in\operatorname{SI}(\chi)\) by an operation under which
\(\operatorname{SI}(\chi)\) is closed.  Hence \(C_\chi\subseteq
\operatorname{SI}(\chi)\).

\emph{Clause 3.}  Let \(\varphi\in\operatorname{SI}(\chi)\).  The set
\(\operatorname{SI}(\chi)\) contains \(\varphi\) and is closed under both
clauses.  The set \(\operatorname{SI}(\varphi)\) is the least set with these
two properties for \(\varphi\), so
\(\operatorname{SI}(\varphi)\subseteq\operatorname{SI}(\chi)\).  If moreover
\(\psi\in\operatorname{SI}(\varphi)\), then \(\psi\in\operatorname{SI}(\chi)\).
\end{proof}
\end{annexitem}

\begin{annexitem}{annex:p02-w02b-p02-020-literal}{thm:p02-w02b-p02-020-literal}{Literal subformulas are instances}
\begin{proof}
\emph{Clause 1.}  First, \(r[x/x]=r\) for every term \(r\), by induction on
the formation of \(r\) in \cref{sec:fol-syntax-semantics-01} using
\cref{sec:fol-syntax-semantics-04}: \(x[x/x]=x\); \(y[x/x]=y\) for a
variable \(y\ne x\); and
\(f(r_1,\ldots,r_n)[x/x]=f(r_1[x/x],\ldots,r_n[x/x])=f(r_1,\ldots,r_n)\) by
the induction hypothesis.  Now prove \(\varphi[x/x]=\varphi\) by strong
induction on the complexity \(c(\varphi)\) of
\cref{sec:fol-syntax-semantics-03}.  For an atomic formula \(r=s\) or
\(R(t_1,\ldots,t_n)\), substitution acts on each term, and each term is
unchanged.  The formula \(\bot\) is unchanged.  For \(\neg\psi\) and
\(\psi\circ\theta\), substitution commutes with the connective, and the
induction hypothesis applies to the immediate subformulas, which have smaller
complexity.  For \(Qy\,\psi\): if \(y=x\), the first clause of
\cref{def:p02-fol-capture-avoiding-substitution} returns \(Qy\,\psi\).  If
\(y\ne x\), then \(y\notin\operatorname{Var}(x)=\{x\}\), so the second clause
applies and returns \(Qy\,(\psi[x/x])\), which is \(Qy\,\psi\) by the
induction hypothesis, since \(c(\psi)<c(Qy\,\psi)\).  The third clause would
require \(y\ne x\) and \(y\in\operatorname{Var}(x)\), which is impossible.

\emph{Clause 2.}  The set \(\operatorname{SI}(\chi)\) contains \(\chi\) and is
closed under passing to immediate propositional subformulas.  It is also closed
under passing from \(Qx\,\varphi\) to \(\varphi\): a variable is a term by
\cref{sec:fol-syntax-semantics-01}, so clause~2 of
\cref{def:p02-w02b-p02-020-instance-set} with \(t=x\) puts
\(\varphi[x/x]\) into \(\operatorname{SI}(\chi)\), and
\(\varphi[x/x]=\varphi\) by clause~1.  The literal subformulas of \(\chi\)
form the least set containing \(\chi\) with these two closure properties, so
they all lie in \(\operatorname{SI}(\chi)\).

\emph{Clause 3.}  Let \(\chi\) be quantifier-free, and let \(P_\chi\) be the
least set of formulas that contains \(\chi\) and is closed under passing to
immediate propositional subformulas.  If \(\neg\psi\) or \(\psi\circ\theta\)
is quantifier-free, so are \(\psi\) and \(\theta\), because every symbol of
\(\psi\) and of \(\theta\) is a symbol of the compound formula.  Hence the set
of quantifier-free members of \(P_\chi\) contains \(\chi\) and has the closure
property, so it is all of \(P_\chi\): every member of \(P_\chi\) is
quantifier-free.  In particular \(P_\chi\) contains no formula of the form
\(Qx\,\theta\), so it is closed, vacuously, under passing from \(Qx\,\theta\)
to \(\theta[t/x]\) and under passing from \(Qx\,\theta\) to \(\theta\).  By
leastness, \(\operatorname{SI}(\chi)\subseteq P_\chi\) and the set of literal
subformulas of \(\chi\) is included in \(P_\chi\).  Conversely, both
\(\operatorname{SI}(\chi)\) and the set of literal subformulas of \(\chi\)
contain \(\chi\) and are closed under passing to immediate propositional
subformulas, so each includes \(P_\chi\).  Therefore
\(\operatorname{SI}(\chi)=P_\chi\), which is the set of literal subformulas of
\(\chi\), and every member is quantifier-free.
\end{proof}
\end{annexitem}

\begin{annexitem}{annex:p02-w02b-p02-020-traced}{thm:p02-w02b-p02-020-traced}{Traced instance property}
\begin{proof}
Say that a pair of formulas \((\psi,\psi')\) is an \emph{admissible step}
when \(\psi=\psi'\), or \(\psi\) is an immediate propositional subformula of
\(\psi'\), or \(\psi'=Qx\,\theta\) and \(\psi=\theta[t/x]\) for some
\(Q\in\{\forall,\exists\}\) and \(t\in\operatorname{Term}(\Lang)\).  These are
the three alternatives of clause~2 of the theorem.  For an admissible step,
\(\psi\in\operatorname{SI}(\psi')\): in the first case because
\(\psi'\in\operatorname{SI}(\psi')\), and in the other two because
\(\operatorname{SI}(\psi')\) contains \(\psi'\) and is closed under the two
clauses of \cref{def:p02-w02b-p02-020-instance-set}.

The height of a derivation is \(0\) for an initial sequent, and one more than
the largest height of the immediate subderivations otherwise.  The proof is
by induction on the height of \(\mathcal D\), for all three clauses
simultaneously.  Let \(S\) be the endsequent of \(\mathcal D\).

\emph{Base: height \(0\).}  Then \(\mathcal D\) is an instance of one of the
two contexted initial sequents of
\texttt{sheet:lk}: \(\mathrm{id}\), the sequent
\(\Gamma,\varphi\Rightarrow\varphi,\Delta\), or \(\bot L\), the sequent
\(\Gamma,\bot\Rightarrow\Delta\), where \(\Gamma,\Delta\) are arbitrary
finite multisets of formulas.  Every formula occurrence \(o\) of
\(\mathcal D\) is an occurrence of \(S\), including the occurrences in
\(\Gamma\) and \(\Delta\), whatever their form.  Its successor chain is
\(o\) alone, with \(k=0\); clause~1 holds, clause~2 is vacuous, and clause~3
holds because \(\psi_0\in\operatorname{SI}(\psi_0)\).

\emph{Inductive step.}  Let \(\mathcal D\) have positive height.  Its last
inference is an instance \(\rho\) of a rule of
\texttt{sheet:lk} with at least one premise, since
\(\mathrm{id}\) and \(\bot L\) occur only as derivations of height \(0\).  By
hypothesis \(\rho\) is neither Cut nor Mix, so it is one of \(W_L\), \(W_R\),
\(C_L\), \(C_R\), \(\neg L\), \(\neg R\), \(\land L\), \(\land R\),
\(\lor L\), \(\lor R\), \(\to L\), \(\to R\), \(\forall L\), \(\forall R\),
\(\exists L\), \(\exists R\); these, with the two initial sequents, Cut and
Mix, are all the rules of the sheet.  Let \(\mathcal D_1\) and, for
\(\land R\), \(\lor L\) and \(\to L\), also \(\mathcal D_2\) be the immediate
subderivations, with endsequents \(S_1\) and \(S_2\), the premises of
\(\rho\).  Each has smaller height and contains neither Cut nor Mix.

Let \(o\) be a formula occurrence of \(\mathcal D\).  If \(o\) lies in \(S\),
its chain is \(o\) alone and the base argument applies.  Otherwise \(o\) lies
in some \(\mathcal D_j\).  The successor of each occurrence of
\(\mathcal D_j\) not in \(S_j\) is fixed by an inference of \(\mathcal D_j\),
so it is the same in \(\mathcal D_j\) and in \(\mathcal D\).  By the
induction hypothesis for \(\mathcal D_j\), the chain of \(o\) in
\(\mathcal D_j\) is \(o=o_0,\ldots,o_m\), with \(o_m\) in \(S_j\), each pair
\((\psi_i,\psi_{i+1})\) with \(i<m\) admissible, and
\(\psi_0\in\operatorname{SI}(\psi_m)\).  In \(\mathcal D\) the occurrence
\(o_m\) is a premise occurrence of \(\rho\), whose successor \(o_{m+1}\) lies
in \(S\) by \cref{def:p02-w02b-p02-020-successor}.  So the chain of \(o\) in
\(\mathcal D\) is \(o_0,\ldots,o_m,o_{m+1}\) with \(k=m+1\), and it ends in
\(S\); this is clause~1.  Once the pair \((\psi_m,\psi_{m+1})\) is shown to be
admissible, with the term named in clause~2, clause~2 holds for the whole
chain.  Clause~3 then follows: \(\psi_m\in\operatorname{SI}(\psi_{m+1})\) by
the remark on admissible steps, and \(\psi_0\in\operatorname{SI}(\psi_m)\), so
\(\psi_0\in\operatorname{SI}(\psi_{m+1})\) by clause~3 of
\cref{thm:p02-w02b-p02-020-closure}.

It remains to check, for each rule, that every premise occurrence \(o'\) of
\(\rho\) and its successor \(s(o')\) carry an admissible pair.  In every rule,
an occurrence \(o'\) in a context multiset \(\Gamma\) or \(\Delta\) of a
premise has as successor the matched occurrence of the same formula in the
conclusion, so the pair is admissible by its first alternative.  This settles
the context occurrences in each case below, including both copies of the
context in the two-premise rules.  The displayed premise formulas are as
follows.
\begin{itemize}
  \item \(W_L\), \(W_R\): the premise \(\Gamma\Rightarrow\Delta\) has only
        context occurrences.  The weakening formula of the conclusion has no
        premise occurrence and needs no check here; it is an occurrence of
        \(S\).
  \item \(C_L\): the premise is \(\Gamma,\varphi,\varphi\Rightarrow\Delta\).
        Each displayed \(\varphi\) has as successor the displayed \(\varphi\)
        of \(\Gamma,\varphi\Rightarrow\Delta\); the pair
        \((\varphi,\varphi)\) is admissible by its first alternative.  \(C_R\)
        is the same on the right.
  \item \(\neg L\): the premise \(\Gamma\Rightarrow\Delta,\varphi\) has
        displayed \(\varphi\), whose successor is \(\neg\varphi\) in
        \(\Gamma,\neg\varphi\Rightarrow\Delta\); \(\varphi\) is the immediate
        propositional subformula of \(\neg\varphi\).  \(\neg R\): the
        displayed \(\varphi\) of \(\Gamma,\varphi\Rightarrow\Delta\) has
        successor \(\neg\varphi\) in \(\Gamma\Rightarrow\Delta,\neg\varphi\),
        by the same alternative.
  \item \(\land L\): the displayed \(\varphi\) and \(\psi\) of
        \(\Gamma,\varphi,\psi\Rightarrow\Delta\) both have successor
        \(\varphi\land\psi\), of which each is an immediate propositional
        subformula.  \(\land R\): the displayed \(\varphi\) of the premise
        \(\Gamma\Rightarrow\Delta,\varphi\) and the displayed \(\psi\) of the
        premise \(\Gamma\Rightarrow\Delta,\psi\) both have successor
        \(\varphi\land\psi\); same alternative.
  \item \(\lor L\): the displayed \(\varphi\) of
        \(\Gamma,\varphi\Rightarrow\Delta\) and the displayed \(\psi\) of
        \(\Gamma,\psi\Rightarrow\Delta\) have successor \(\varphi\lor\psi\).
        \(\lor R\): the displayed \(\varphi\) and \(\psi\) of
        \(\Gamma\Rightarrow\Delta,\varphi,\psi\) have successor
        \(\varphi\lor\psi\).  In both, each is an immediate propositional
        subformula of \(\varphi\lor\psi\).
  \item \(\to L\): the displayed \(\varphi\) of the premise
        \(\Gamma\Rightarrow\Delta,\varphi\) and the displayed \(\psi\) of the
        premise \(\Gamma,\psi\Rightarrow\Delta\) have successor
        \(\varphi\to\psi\) in \(\Gamma,\varphi\to\psi\Rightarrow\Delta\).
        \(\to R\): the displayed antecedent \(\varphi\) and succedent \(\psi\)
        of \(\Gamma,\varphi\Rightarrow\Delta,\psi\) have successor
        \(\varphi\to\psi\) in \(\Gamma\Rightarrow\Delta,\varphi\to\psi\).  In
        both, each is an immediate propositional subformula of
        \(\varphi\to\psi\).
  \item \(\forall L\): the displayed \(\varphi[t/x]\) of
        \(\Gamma,\varphi[t/x]\Rightarrow\Delta\) has successor
        \(\forall x\,\varphi\).  By the side condition of the rule,
        \([t/x]\) is the operation of
        \cref{def:p02-fol-capture-avoiding-substitution}, the same operation
        as in clause~2 of \cref{def:p02-w02b-p02-020-instance-set}, and \(t\)
        is a term of \(\Lang\).  The pair is admissible by its third
        alternative, with \(Q=\forall\), \(\theta=\varphi\) and the term
        \(t\) of the rule instance.  \(\exists R\): the displayed
        \(\varphi[t/x]\) of \(\Gamma\Rightarrow\Delta,\varphi[t/x]\) has
        successor \(\exists x\,\varphi\); the same alternative with
        \(Q=\exists\).
  \item \(\forall R\): the displayed \(\varphi[y/x]\) of
        \(\Gamma\Rightarrow\Delta,\varphi[y/x]\) has successor
        \(\forall x\,\varphi\).  The eigenvariable \(y\) is a variable, hence
        a term of \(\Lang\) by \cref{sec:fol-syntax-semantics-01}, and the substitution is
        again the operation of
        \cref{def:p02-fol-capture-avoiding-substitution}.  The pair is
        admissible by its third alternative with \(Q=\forall\),
        \(\theta=\varphi\) and \(t=y\).  \(\exists L\): the displayed
        \(\varphi[y/x]\) of \(\Gamma,\varphi[y/x]\Rightarrow\Delta\) has
        successor \(\exists x\,\varphi\); the same alternative with
        \(Q=\exists\) and \(t=y\).  In both rules the freshness condition on
        \(y\) holds because \(\mathcal D\) is a derivation on the sheet; the
        verification of admissibility does not depend on it.
\end{itemize}
Every rule other than Cut and Mix has been treated, which completes the
inductive step.

For the final assertion, let \(o\) be any formula occurrence of
\(\mathcal D\), and let \(\chi\) be the formula at the last occurrence
\(o_k\) of its chain, an occurrence of the endsequent.  By clause~3 the
formula at \(o\) lies in \(\operatorname{SI}(\chi)\), so by clause~1 of
\cref{thm:p02-w02b-p02-020-closure} it is a first-order subformula instance of
\(\chi\) in the sense of \cref{def:p02-fol-subformula-instance}.  The steps of
the chain that pass through \(\forall L\) or \(\exists R\) use the term
instances of those rules, and those that pass through \(\forall R\) or
\(\exists L\) use the eigenvariable instances, as recorded in clause~2.
\end{proof}
\end{annexitem}

\begin{annexitem}{annex:p02-w02b-p02-020-quantifier-free}{thm:p02-w02b-p02-020-quantifier-free}{Quantifier-free endsequents give literal subformulas}
\begin{proof}
Let \(o\) be a formula occurrence of \(\mathcal D\) with formula \(\psi\).  By
\cref{thm:p02-w02b-p02-020-traced}, \(\psi\in\operatorname{SI}(\chi)\) for the
formula \(\chi\) at some occurrence of the endsequent.  By hypothesis
\(\chi\) is quantifier-free, so by clause~3 of
\cref{thm:p02-w02b-p02-020-literal}, \(\operatorname{SI}(\chi)\) is the set of
literal subformulas of \(\chi\) and all its members are quantifier-free.
Hence \(\psi\) is a quantifier-free literal subformula of \(\chi\).  If an
instance of \(\forall L\), \(\forall R\), \(\exists L\) or \(\exists R\)
occurred in \(\mathcal D\), its principal formula \(Qx\,\varphi\) would be a
formula occurring in \(\mathcal D\) that contains a quantifier, which has just
been excluded.
\end{proof}
\end{annexitem}

\begin{annexitem}{annex:p02-w02b-p02-020-host-property}{thm:p02-fol-subformula-property}{Subformula property for cut-free first-order \(\mathsf{LK}\)}
\begin{proof}
Let \(\mathcal D\) be a derivation on \texttt{sheet:lk}
using neither Cut nor Mix, and let \(o\) be a formula occurrence of
\(\mathcal D\).  By \cref{thm:p02-w02b-p02-020-traced}, the successor chain of
\(o\) ends at an occurrence of the endsequent with formula \(\chi\), and the
formula at \(o\) lies in \(\operatorname{SI}(\chi)\); by clause~1 of
\cref{thm:p02-w02b-p02-020-closure} it is a first-order subformula instance of
\(\chi\) in the sense of \cref{def:p02-fol-subformula-instance}.  The steps of
the chain made by \(\forall L\) and \(\exists R\) are term instances, and
those made by \(\forall R\) and \(\exists L\) are eigenvariable instances
(clause~2 of \cref{thm:p02-w02b-p02-020-traced}).  The claim is not literal
subformulahood: \cref{def:p02-w02b-p02-020-example} gives a derivation with a
formula that is no literal subformula of the endsequent.  It does not extend to
derivations with Cut or Mix (\cref{def:p02-w02b-p02-020-cut-failure}).
\end{proof}
\end{annexitem}
